\documentclass[11pt]{article}

\usepackage[margin=1in]{geometry}
\usepackage{amsmath}
\usepackage{tcolorbox}
\usepackage{xcolor}
\usepackage{array}
\usepackage{longtable}
\usepackage{enumitem}

\definecolor{secondary}{RGB}{230,240,255}
\definecolor{warning}{RGB}{255,245,220}
\definecolor{insight}{RGB}{235,255,240}

\title{\textbf{NVIDIA Deep Learning Software Engineer Interview Review}\\
Algorithmic Model Optimization (LLMs, Diffusion, TensorRT, TRT-LLM)}
\author{}
\date{}

\begin{document}
\maketitle
\renewcommand{\arraystretch}{1.2}

\section*{Role-Focused Quick Outline}
\begin{itemize}[leftmargin=*]
\item Model optimization: quantization, sparsity, pruning, distillation, NAS
\item LLM inference bottlenecks: prefill vs decode, KV cache, batching
\item PyTorch 2.0 stack: TorchDynamo, torch.export, torch.compile, FX/IR
\item CUDA/TRT-LLM/Triton kernel-level performance optimization
\item Multi-GPU inference: tensor, pipeline, sequence parallelism
\item Profiling and debugging: Nsight Systems, Nsight Compute, PyTorch profiler
\item Deployment platform thinking: automation, validation, rollback, observability
\item High-probability interview Q\&A and final review cheat sheet
\end{itemize}

\section{Model Optimization Techniques}

\subsection{Quantization (Highest Probability)}
\begin{tcolorbox}[colback=secondary]
\textbf{Core insight:} Quantization usually helps latency by reducing memory traffic first; raw compute savings are secondary unless kernels are optimized end-to-end.
\end{tcolorbox}

\subsubsection*{PTQ vs QAT}
\begin{center}
\begin{tabular}{|p{0.26\linewidth}|p{0.33\linewidth}|p{0.33\linewidth}|}
\hline
\textbf{Method} & \textbf{Advantages} & \textbf{Tradeoffs} \\
\hline
PTQ (Post-Training Quantization) & Fast to apply, low engineering overhead & Higher accuracy risk, calibration quality matters \\
\hline
QAT (Quantization-Aware Training) & Better accuracy retention at low bit widths & Requires retraining and more MLOps effort \\
\hline
\end{tabular}
\end{center}

\textbf{Why quantization may fail to improve latency:}
\begin{itemize}[leftmargin=*]
\item Unsupported INT8/FP8 kernels force fallback to FP16/FP32
\item Dequantize/requantize boundaries add memory and launch overhead
\item Fusion opportunities are lost due to mixed dtypes
\item Workload is already compute-saturated at target batch size
\item Poor calibration causes unstable activations and recovery logic
\end{itemize}

\subsection{Sparsity and Pruning}
\begin{center}
\begin{tabular}{|p{0.22\linewidth}|p{0.35\linewidth}|p{0.35\linewidth}|}
\hline
\textbf{Type} & \textbf{Description} & \textbf{Practical impact} \\
\hline
Unstructured sparsity & Removes individual weights & Often limited speedup unless specialized sparse kernels exist \\
\hline
Structured sparsity & Removes heads/channels/blocks & Better hardware mapping and real latency wins \\
\hline
\end{tabular}
\end{center}

\begin{tcolorbox}[colback=insight]
\textbf{Interview signal:} ``Sparsity is only useful when compiler + kernel + hardware path is sparse-aware.''
\end{tcolorbox}

\subsection{Distillation and NAS}
\begin{itemize}[leftmargin=*]
\item Distillation compresses quality into smaller/faster student models
\item NAS can discover hardware-friendly blocks but has high search cost
\item For production, emphasize \textbf{search cost vs deployment ROI}
\end{itemize}

\section{LLM Inference Optimization}

\subsection{Prefill vs Decode}
\begin{center}
\begin{tabular}{|p{0.2\linewidth}|p{0.28\linewidth}|p{0.44\linewidth}|}
\hline
\textbf{Phase} & \textbf{Primary bottleneck} & \textbf{Best optimization levers} \\
\hline
Prefill & GEMM compute + activation memory & Better batching, tensor core utilization, graph capture, fusion \\
\hline
Decode & KV-cache bandwidth/latency & Paged KV cache, attention kernel fusion, reduced precision cache, continuous batching \\
\hline
\end{tabular}
\end{center}

\subsection{Parallelism and Sharding}
\begin{itemize}[leftmargin=*]
\item Tensor parallelism: split matrix dimensions across GPUs
\item Pipeline parallelism: split layers across stages
\item Sequence parallelism: split sequence/tokens across devices
\item Data parallel inference: replicate model and route requests
\end{itemize}

\textbf{Scaling blockers:}
\begin{itemize}[leftmargin=*]
\item NCCL communication overhead dominates after a point
\item Interconnect mismatch (PCIe vs NVLink vs InfiniBand)
\item Micro-batch too small to hide communication
\item Host-side scheduling and data pipeline stalls
\end{itemize}

\section{GPU Architecture and Profiling}

\subsection{Performance Limits}
\begin{center}
\begin{tabular}{|p{0.3\linewidth}|p{0.62\linewidth}|}
\hline
\textbf{Constraint} & \textbf{Typical cause in LLM inference} \\
\hline
Memory bandwidth & Large weights + KV-cache streaming each decode step \\
\hline
Compute throughput & Dense GEMMs and attention math in prefill \\
\hline
Kernel launch overhead & Too many tiny kernels or poor fusion \\
\hline
Low occupancy & Register/shared-memory pressure limits active warps \\
\hline
\end{tabular}
\end{center}

\subsection{Essential Tools}
\begin{center}
\begin{tabular}{|p{0.3\linewidth}|p{0.62\linewidth}|}
\hline
\textbf{Tool} & \textbf{What you use it for} \\
\hline
Nsight Systems & End-to-end timeline, launch gaps, communication overlap \\
\hline
Nsight Compute & Kernel-level metrics: occupancy, memory transactions, tensor core usage \\
\hline
PyTorch Profiler & Operator hotspots, shape-dependent regressions, eager vs compiled comparison \\
\hline
nvidia-smi/DCGM & Device health, utilization, memory pressure, throttling signals \\
\hline
\end{tabular}
\end{center}

\textbf{Benchmark hygiene checklist:}
\begin{itemize}[leftmargin=*]
\item Warmup iterations before measurement
\item Synchronize correctly for GPU timing
\item Report p50/p95 and throughput together
\item Fix prompt length, output length, and batch policy
\item Separate prefill and decode metrics
\end{itemize}

\section{PyTorch 2.0 to Deployment Stack}
\begin{itemize}[leftmargin=*]
\item TorchDynamo captures graphs from eager Python programs
\item \texttt{torch.export} creates a normalized, portable graph representation
\item \texttt{torch.compile} applies backend optimizations and fusion
\item Exported graphs feed TensorRT/TRT-LLM style automated deployment pipelines
\end{itemize}

\begin{tcolorbox}[colback=secondary]
\textbf{Great interview framing:} ``I use graph capture to reduce Python overhead, unlock fusion, and produce stable artifacts for automated deployment and regression testing.''
\end{tcolorbox}

\section{High-Probability Interview Q\&A (Role Specific)}
\begin{center}
\begin{tabular}{|p{0.34\linewidth}|p{0.58\linewidth}|}
\hline
\textbf{Question} & \textbf{Strong answer (concise)} \\
\hline
How would you optimize LLM inference end-to-end? & Baseline first, split prefill/decode metrics, find top kernels with Nsight, optimize memory path (KV cache + fusion), then validate quality and rollback safety. \\
\hline
Why does GPU utilization look low during serving? & Decode is often memory-bound and request-driven; low SM utilization can still be optimal if latency is constrained. Check memory BW, kernel launch gaps, and batching policy. \\
\hline
When does quantization hurt? & If kernels fall back, if dtype boundaries add conversion overhead, or if quality drops force rework. I verify kernel coverage and quality gates before rollout. \\
\hline
TorchDynamo vs torch.export? & Dynamo captures runtime graphs from Python execution; export gives a cleaner, standardized graph for stable deployment transformations. \\
\hline
How do you debug scaling stalls after adding GPUs? & Measure compute/comm overlap, inspect NCCL traces, verify topology and message sizes, and adjust parallelism mode and micro-batch size. \\
\hline
What makes attention kernels fast? & Coalesced memory access, fused ops, low overhead KV-cache reads/writes, and good tensor core utilization for compatible shapes/dtypes. \\
\hline
How do you choose tensor vs pipeline parallelism? & Tensor parallel for large GEMMs and balanced layers; pipeline for very deep models or memory limits. Choice depends on interconnect and latency target. \\
\hline
How do you ensure safe deployment of optimizations? & Keep golden baselines, automated accuracy/perf gates, versioned artifacts, canary rollout, and fast rollback if SLOs or quality drift regress. \\
\hline
What would you optimize in TRT/TRT-LLM integration? & Graph lowering coverage, kernel selection heuristics, KV-cache policies, and reproducible benchmark harnesses across model families. \\
\hline
How do you reason about roofline for transformers? & Decode tends toward memory roof due to KV traffic; prefill can approach compute roof with larger batches and fused GEMM-heavy paths. \\
\hline
How do you improve user-perceived latency? & Continuous batching, token streaming, smart scheduling, and prioritizing p95 TTFT/token latency over only peak throughput. \\
\hline
What tests are mandatory for optimization code? & Numerical parity tolerances, shape/dtype coverage, regression benchmarks, and stress tests for long-context KV-cache behavior. \\
\hline
\end{tabular}
\end{center}

\section{System Design Prompts You Should Practice}
\begin{enumerate}[leftmargin=*]
\item Design an automated model optimization service from PyTorch model input to TensorRT/TRT-LLM deployment artifact.
\item Design multi-tenant LLM serving with dynamic batching, per-request QoS, and safe rollout.
\item Design a profiling pipeline that automatically flags kernel regressions after code changes.
\item Design model sharding strategy for long-context inference across 8+ GPUs.
\end{enumerate}

\section{Software Engineering Signals for This Role}
\begin{itemize}[leftmargin=*]
\item Strong debugging and performance analysis loop
\item Modular interfaces across model, compiler, kernel, runtime layers
\item Test design that catches numerical and performance regressions
\item Clear documentation and reproducible benchmark methodology
\end{itemize}

\section{Final Table Cheat Sheet: What to Review}
\begin{longtable}{|p{0.16\linewidth}|p{0.09\linewidth}|p{0.24\linewidth}|p{0.31\linewidth}|}
\hline
\textbf{Topic} & \textbf{Priority} & \textbf{What to be ready to explain} & \textbf{Detailed explanation to prepare} \\
\hline
\endfirsthead
\hline
\textbf{Topic} & \textbf{Priority} & \textbf{What to be ready to explain} & \textbf{Detailed explanation to prepare} \\
\hline
\endhead
\hline
\endfoot
\hline
\endlastfoot
Quantization (INT8/FP8, PTQ/QAT) & High & Accuracy tradeoffs, kernel coverage, why latency may not improve & Explain that speedup comes only when the full path is quantized: graph lowering, kernels, and memory format. Be ready to discuss calibration datasets, sensitive ops (LayerNorm/Softmax), and fallback overhead from dequantize/requantize boundaries. \\
\hline
KV cache and attention optimization & High & Paged KV, decode bottlenecks, memory layout and bandwidth implications & Decode is usually memory-bound because each new token reads prior K/V states. Cover paged KV cache, contiguous block layout, cache precision choices, and fused attention kernels to reduce memory traffic and kernel launch overhead. \\
\hline
PyTorch 2.0 compiler stack & High & TorchDynamo, torch.export, torch.compile, graph capture value & Clarify differences: TorchDynamo captures runtime graphs, \texttt{torch.export} creates stable deployable graphs, and \texttt{torch.compile} applies backend optimizations/fusion. Explain why this reduces Python overhead and improves portability to deployment backends. \\
\hline
TensorRT/TRT-LLM basics & High & Deployment path, engine building, optimization and benchmarking workflow & Walk through model export, graph transformation, engine build/tactic selection, runtime execution, and regression validation. Mention dynamic shapes, precision configuration, plugin coverage, and performance-quality checks before rollout. \\
\hline
CUDA/Triton kernel fundamentals & High & Occupancy, memory coalescing, fusion, tensor core-friendly tiling & Be ready to reason from bottleneck to kernel design: choose tile sizes, manage shared memory/register pressure, maximize coalesced loads, and map compute to tensor cores. Explain tradeoffs between occupancy and per-thread resource usage. \\
\hline
Distributed inference parallelism & High & Tensor/pipeline/sequence parallelism and communication bottlenecks & Explain when to choose each parallelism mode based on model size, latency target, and topology. Discuss all-reduce/all-gather costs, overlap strategies, micro-batch tuning, and why scaling can plateau when communication dominates compute. \\
\hline
Profiling toolchain & High & Nsight Systems vs Nsight Compute vs PyTorch profiler usage & Present a profiling workflow: use Nsight Systems for end-to-end timeline and stalls, Nsight Compute for kernel metrics (occupancy, memory transactions, tensor core utilization), and PyTorch profiler for operator-level hotspots tied to model code. \\
\hline
GPU architecture essentials & Medium & SMs, warps, memory hierarchy, compute vs bandwidth limits & Explain how warp scheduling hides latency, how memory hierarchy affects throughput, and why roofline thinking helps classify kernels as compute- or memory-bound. Connect this directly to transformer prefill (more compute-heavy) vs decode (more bandwidth-bound). \\
\hline
Benchmark methodology & High & Warmup, synchronization, p50/p95, TTFT, tokens/sec, fair comparisons & Be precise on measurement hygiene: fixed prompt/output lengths, warmup passes, synchronized timing, percentile reporting, and separate prefill/decode analysis. Explain that fair comparisons require equal batching policy, model config, and precision settings. \\
\hline
Reliability and deployment safety & Medium & Canary rollout, rollback, artifact versioning, validation gates & Describe production safeguards: versioned optimization artifacts, automated numerical/performance gates, canary release with SLO monitoring, and deterministic rollback. Include how to track drift and regressions across model and runtime updates. \\
\hline
\end{longtable}

\begin{tcolorbox}[colback=insight]
\textbf{Final takeaway:} For this NVIDIA role, your strongest interview narrative is: \textit{profile first, optimize where bottlenecks are real, and ship safe automated deployment at scale.}
\end{tcolorbox}

\end{document}
